
In this section we explain the finite inputs and scalar estimates from
Ref.~\cite{OpenAI2026Periodic} that certify the starting pair
\begin{align}
 1-\Ps(36,49/4)&\le\frac{47916}{10^6},\notag \\
 1-\Pt(72,49/4)&\le\frac{181521}{10^6}.
 \label{cert:eq:base}
\end{align}
These two bounds form Lemma~\ref{lem:arxiv-seed36} in the main text.
We then present the temperature-lifting calculation, which combines the
same short-ring trace bounds with the $120$-site trial state of
Ref.~\cite{OpenAI2026Periodic} to prove a unique ground state and
\begin{equation}
 \Delta_L>0.072901771
 \quad\text{for every even }L\ge36.
 \label{cert:eq:improved}
\end{equation}
We have reproduced the finite trace and trial contractions of
Ref.~\cite{OpenAI2026Periodic} in exact arithmetic; the stronger gap bound
uses only additional scalar inequalities on those inputs.
Every terminating decimal that defines a bound below denotes an exact
rational number, and we mark numerical approximations explicitly.

\subsection{The base seed at length 36}
\label{cert:sec:base}

To prove Eq.~\eqref{cert:eq:base} at $b_0=49/4$, we need lower bounds on
the numerators $Z_{72}(b_0)$ and $Z_{72}(2b_0)$ and upper bounds on the
denominators $Z_{36}(b_0)^2$ and $Z_{72}(b_0)^2$. A $72$-site trial state
bounds $Z_{72}(b)$ from below for every $b>0$. Traces on $4\le n\le12$
sites constrain the eigenvalues of the one-site spatial transfer operator
at $b_0$, which lets us bound the subleading contributions to
$Z_{36}(b_0)$ and $Z_{72}(b_0)$.

\subsubsection{Hamiltonians, twists, and the spatial transfer operator}
\label{cert:sec:transfer}

We consider the periodic spin-1 Heisenberg Hamiltonian
\begin{equation}
 H_n=\sum_{j=1}^n\bm S_j\cdot\bm S_{j+1},
 \qquad \bm S_j^2=2I,\qquad \bm S_{n+1}=\bm S_1,
 \label{cert:eq:hamiltonian}
\end{equation}
with unit exchange coupling. Lengths count sites, and
$\Delta_L=E_1(L)-E_0(L)$, with eigenvalues ordered with multiplicity.
Anticipating the trial energy per site, we shift the Hamiltonian by
$\kappa$ per site:
\begin{equation}
 \kappa=\frac{700741}{500000},\qquad
 \widehat H_n=H_n+\kappa nI,\qquad
 Z_n(\beta)=\Tr e^{-\beta\widehat H_n}.
 \label{cert:eq:shift}
\end{equation}
This extensive shift cancels in both purities:
\begin{equation}
 \Ps(n,\beta)=\frac{Z_{2n}(\beta)}{Z_n(\beta)^2},
 \qquad
 \Pt(n,\beta)=\frac{Z_n(2\beta)}{Z_n(\beta)^2}.
 \label{cert:eq:purities}
\end{equation}

To separate the spatial transfer spectrum by spin-rotation symmetry, we
also use rings with a twisted closing bond. For an on-site spin rotation
$u_g$ representing $g$, we define
\begin{align}
 H_n(g)&=\sum_{j=1}^{n-1}h_{j,j+1}+h_{n,1}(g),\notag \\
 h_{n,1}(g)&=\sum_{\alpha=x,y,z}
 S_n^\alpha(u_gS^\alpha u_g^\dagger)_1,\notag \\
 Z_n(\beta,g)&=\Tr e^{-\beta(H_n(g)+\kappa nI)},
 \label{cert:eq:twisted-hamiltonian}
\end{align}
where $h_{j,j+1}=\bm S_j\cdot\bm S_{j+1}$.
The twist modifies the spatial seam rather than inserting $u_g^{\otimes n}$
into the untwisted Gibbs trace.

The continuous-time construction of Ref.~\cite{OpenAI2026Periodic}
provides a self-adjoint one-site spatial transfer operator $K_\beta$ for
the shifted Hamiltonian and a virtual symmetry operator $U_g$ such that
\begin{equation}
 Z_n(\beta,g)=\Tr(K_\beta^nU_g),\qquad n\ge4.
 \label{cert:eq:transfer-trace}
\end{equation}
At $g=1$, $Z_n(\beta)=\sum_i\lambda_i(\beta)^n$ is a power sum of the
real eigenvalues $\lambda_i(\beta)$ of $K_\beta$. Because the same
eigenvalues govern every length $n\ge4$, trace bounds on $4\le n\le12$
constrain the power sums at $n=36$ and $n=72$. Unlike the positive
two-site Trotter matrix of Section~\ref{sm:sec:transfer}, the one-site
operator $K_\beta$ can have negative eigenvalues, so we must bound
$|\lambda_i(\beta)|$ on both sides of zero.

In our diagrams, horizontal wires carry spatial virtual indices, vertical
double wires carry physical indices, and closed loops denote traces.
Commutation of $K_\beta$ with $U_g$ gives
\begin{align}
 \tnchain[last=sym]{K_\beta,U_g}
 &=\tnchain[node=sym,last=xop]{U_g,K_\beta},\notag \\
 Z_n(\beta,g)
 &=\tnchain[node=xvertex,last=sym,ring]{,,\cdots,,U_g},
 \label{cert:eq:twist-diagram}
\end{align}
where the ring contains $n$ transfer vertices and $U_g$ implements the
twisted closing bond of Eq.~\eqref{cert:eq:twisted-hamiltonian}.

\subsubsection{Exact finite trace and trial inputs}
\label{cert:sec:inputs}

We now state the two numerical inputs from Ref.~\cite{OpenAI2026Periodic}
used at $b_0=49/4$. The first input is the table of shifted traces
$Z_n(b_0)$ and $Z_n(b_0,P)$ for $4\le n\le12$, where $P$ is a $\pi$
rotation about a Cartesian axis. Table~\ref{cert:tab:traces} lists the
rational centers, each certified within the absolute error
\begin{equation}
 \epsilon_Z=3\times10^{-5}.
 \label{cert:eq:trace-error}
\end{equation}

\begin{table}[t]
 \caption{Exact rational centers of the shifted traces at $b_0=49/4$ from
 Ref.~\cite{OpenAI2026Periodic}. Each entry has certified absolute error
 at most $3\times10^{-5}$.}
 \label{cert:tab:traces}
 \begin{ruledtabular}
 \begin{tabular}{r rr}
 $n$ & $Z_n(b_0)$ center & $Z_n(b_0,P)$ center\\
 \hline
 4  & 124.885379 & 0.065757\\
 5  & 0.002925   & 2.960742\\
 6  & 12.870614  & 0.286599\\
 7  & 0.054069   & 1.920414\\
 8  & 4.639283   & 0.510140\\
 9  & 0.195455   & 1.514846\\
 10 & 2.654823   & 0.675419\\
 11 & 0.377629   & 1.314544\\
 12 & 1.900547   & 0.787405
 \end{tabular}
 \end{ruledtabular}
\end{table}

These enclosures come from integer Horner recursions for the half-Gibbs
operator rather than dense diagonalization. For an unshifted Hamiltonian
$H=H_n(g)$, we set
\begin{equation}
 R=I-\frac{(3n/2)I+H}{16},
 \label{cert:eq:uniformization}
\end{equation}
which satisfies $\|R\|\le1$ on each ring in Table~\ref{cert:tab:traces}.
At $b_0=49/4$, the shifted half-Gibbs operator is
\begin{equation}
 e^{-b_0(H+\kappa nI)/2}
 =e^{-b_0(\kappa-3/2)n/2}e^{98(R-I)}.
 \label{cert:eq:half-gibbs}
\end{equation}
We approximate $e^{98(R-I)}$ by the normalized Poisson polynomial
\begin{equation}
 \sum_{j=0}^{J}p_jR^j,\qquad
 p_j=\frac{98^j/j!}{\sum_{\ell=0}^{J}98^\ell/\ell!},
 \qquad J=280,
 \label{cert:eq:poisson}
\end{equation}
whose truncation and normalization error in operator norm is at most
\begin{equation}
 2e^{-98}\sum_{j>280}\frac{98^j}{j!}<2^{-84}.
 \label{cert:eq:poisson-error}
\end{equation}

Setting $Q=2^{55}$, we replace $Qp_j$ by integers $c_j$ with
$|c_j-Qp_j|\le1$. Writing $R=D/s$ with an integer matrix $D$ and integer
divisor $s>0$, we evaluate each basis column $e_r$ by the downward
recursion
\begin{equation}
 X_{J+1}=0,\qquad
 X_j=\left\lfloor\frac{DX_{j+1}}s\right\rfloor+c_je_r,
 \label{cert:eq:horner}
\end{equation}
for $j=J,\ldots,0$, where the floor acts entrywise, including on negative
entries.
In dimension $d=3^n$, each step introduces a Euclidean rounding error of
at most $\sqrt d+1$. Because $\|R\|\le1$, these errors do not amplify, so
the total Euclidean error in $X_0$ is at most $(J+1)(\sqrt d+1)$ before
adding the Poisson tail~\eqref{cert:eq:poisson-error}.
The trace $Z_n(b_0,g)$ equals the squared Frobenius norm of the
half-Gibbs operator, which we obtain by summing the squared column norms
with their exact symmetry multiplicities and multiplying by a rational
enclosure of the scalar prefactor in Eq.~\eqref{cert:eq:half-gibbs}.

The second input is a periodic matrix-product trial state defined by real
integer matrices $A_s$ of bond dimension $26$, used at $N=72$ for the base
seed and at $N=120$ for temperature lifting. Setting
\begin{align}
 \psi_N(s_1,\ldots,s_N)&=\Tr(A_{s_1}\cdots A_{s_N}),\notag \\
 T&=\sum_s A_s\otimes A_s,\notag \\
 O&=\sum_{s,t,s',t'}h_{st,s't'}
       A_sA_t\otimes A_{s'}A_{t'},
 \label{cert:eq:mps}
\end{align}
where $T$ is the $26^2\times26^2$ trial transfer matrix, we express the
norm and energy expectation as exact finite-matrix traces:
\begin{align}
 \|\psi_N\|^2&=\seedmps0=\Tr T^N,\notag \\
 \langle\psi_N,H_N\psi_N\rangle
   &=N\,\seedmps1=N\Tr(T^{N-2}O).
 \label{cert:eq:mps-contractions}
\end{align}
Each layer has $N$ tensors, and the two-site operator $h$ acts on one
bond. Exact integer arithmetic verifies that for $N\in\{72,120\}$,
\begin{align}
 &\Tr T^N>0,\notag \\
 &500000\Tr(T^{N-2}O)+700741\Tr T^N<0.
 \label{cert:eq:trial-input}
\end{align}
By the variational principle, the ground energy satisfies
$E_0(N)<-\kappa N$, so the shifted ground energy $E_0(N)+\kappa N$ is
negative. Retaining only the ground-state Boltzmann weight gives the lower
bounds
\begin{equation}
 Z_{72}(\beta)>1,\qquad Z_{120}(\beta)>1
 \quad(\beta>0).
 \label{cert:eq:trial-consequence}
\end{equation}
In the base seed, $Z_{72}(2b_0)>1$ bounds the thermal numerator directly,
whereas $Z_{72}(b_0)>1$ forces a large leading spatial eigenvalue.

\subsubsection{Sector moments and polynomial filters}
\label{cert:sec:filters}

We next split the spectrum of $K_{b_0}$ into spin-rotation sectors and
bound its individual eigenvalues. The four virtual $\pi$ rotations
$D_2=\{1,P_x,P_y,P_z\}$ commute with $K_{b_0}$ and define the character
projectors
\begin{equation}
 \Pi_\chi=\frac14\sum_{g\in D_2}\chi(g)U_g,\qquad
 \Tr(K_{b_0}^n\Pi_\chi)
 =\tnchain[node=xvertex,last=sym,ring]{,,\cdots,,\Pi_\chi}.
 \label{cert:eq:projectors}
\end{equation}
Cyclic permutation of the coordinate axes conjugates the three nontrivial
characters into one another, so their three sectors have identical
eigenvalue lists, which we denote by $N$; we denote the $D_2$-invariant
eigenvalue list by $J$. Writing $J_n=\sum_{\lambda\in J}\lambda^n$ and
$N_n=\sum_{\lambda\in N}\lambda^n$ at $b_0$, we have
\begin{align}
 Z_n&=J_n+3N_n,& Z_n(P)&=J_n-N_n,\notag \\
 J_n&=\frac{Z_n+3Z_n(P)}4,&
 N_n&=\frac{Z_n-Z_n(P)}4.
 \label{cert:eq:sector-inversion}
\end{align}
From Eq.~\eqref{cert:eq:trace-error}, the sector moments $J_n$ and $N_n$
for $4\le n\le12$ have absolute errors at most $\epsilon_Z$ and
$\epsilon_Z/2$ around their rational centers. In particular, at $n=12$
these enclosures yield the twelfth-moment upper bounds
\begin{equation}
 J_{12}\le m_J=\frac{2131441}{2000000},\qquad
 N_{12}\le m_N=\frac{556631}{2000000}.
 \label{cert:eq:masses}
\end{equation}

To bound individual eigenvalues from these finite power sums, we evaluate
the nonnegative polynomials of Ref.~\cite{OpenAI2026Periodic},
\begin{align}
 F_J(x)&=x^4(7+17x-280x^2-185x^3+1000x^4)^2,\notag \\
 F_N(x)&=x^4(12-394x^2+1000x^4)^2.
 \label{cert:eq:filters}
\end{align}
Pairing each coefficient of $F_J$ and $F_N$ with the upper or lower
endpoint of the corresponding moment interval for $4\le n\le12$ gives
\begin{align}
 \sum_{\lambda\in J}F_J(\lambda)&\le B_J=318604.54848175,\notag \\
 \sum_{\lambda\in N}F_N(\lambda)&\le B_N=48169.880699.
 \label{cert:eq:filter-bounds}
\end{align}
Because $F_J$ and $F_N$ are nonnegative on $\mathbb R$, each single
eigenvalue satisfies $F_J(\lambda)\le B_J$ for $\lambda\in J$ and
$F_N(\lambda)\le B_N$ for $\lambda\in N$.

For $U=1.00139$ and $V=0.872$, we expand $F_J(\pm(U+y))$ and
$F_N(\pm(V+y))$ in powers of $y\ge0$. Exact rational arithmetic shows that
their constant terms exceed $B_J$ and $B_N$ and all higher coefficients
are nonnegative, which rules out eigenvalues on the exterior rays
$|x|\ge U$ and $|x|\ge V$:
\begin{equation}
 |\lambda|<U\quad(\lambda\in J),\qquad
 |\lambda|<V\quad(\lambda\in N).
 \label{cert:eq:caps}
\end{equation}
Checking both signs $\pm(U+y)$ and $\pm(V+y)$ is required because
$K_{b_0}$ is self-adjoint rather than positive.

\subsubsection{Capped masses and the two base defects}
\label{cert:sec:base-defects}

To pass from the twelfth moments~\eqref{cert:eq:masses} and the eigenvalue
caps~\eqref{cert:eq:caps} to $n=36$ and $n=72$, we bound higher power sums
under a mass constraint and an entry cap. For $m\ge0$, $a>0$, and $r\ge1$,
set
\begin{equation}
 \mathcal C(m,a,r)=qa^r+(m-qa)^r,\qquad
 q=\lfloor m/a\rfloor.
 \label{cert:eq:capped}
\end{equation}
Whenever $0\le t_i\le a$ and $\sum_i t_i\le m$, convexity of $t^r$
implies $\sum_i t_i^r\le\mathcal C(m,a,r)$: shifting mass from a smaller
unsaturated entry to a larger one increases $\sum_i t_i^r$, so the maximum
is achieved by placing $q$ entries at the cap $a$ and the remaining mass
$m-qa$ in a single entry. Applied to $t_i=\lambda_i^{12}\ge0$, this bounds
higher moments of the $12$th-power weights.

We first show that $J$ contains a dominant eigenvalue with
$|\lambda_0|>\ell=0.999$. If every $\lambda\in J$ satisfied
$|\lambda|\le\ell$, then applying Eq.~\eqref{cert:eq:capped} with $r=6$
to $J$ (with cap $\ell^{12}$) and to each of the three $N$ sectors (with
cap $V^{12}$) would give
\begin{equation}
 Z_{72}(b_0)
 \le\mathcal C(m_J,\ell^{12},6)
      +3\mathcal C(m_N,V^{12},6)
 <0.931<1,
 \label{cert:eq:dominance-check}
\end{equation}
contradicting the $72$-site trial bound $Z_{72}(b_0)>1$ in
Eq.~\eqref{cert:eq:trial-consequence}. Because $m_N<\ell^{12}$, no $N$
eigenvalue can exceed $\ell$ in modulus, so $J$ has at least one
eigenvalue $\lambda_0$ with $|\lambda_0|>\ell$.

Removing $\lambda_0^{12}>\ell^{12}$ from $J_{12}$ leaves residual twelfth
moment at most $m_J-\ell^{12}<U^{12}$ in $J$. For even $k\ge12$, the
subleading moment $\sum_{i\ne0}\lambda_i^k$ is therefore bounded by
\begin{equation}
 R_k=(m_J-\ell^{12})^{k/12}
       +3\mathcal C(m_N,V^{12},k/12).
 \label{cert:eq:base-residual}
\end{equation}
At $k=36$ and $k=72$, we obtain the two base purity bounds:
\begin{align}
 \Ps(36,b_0)&\ge
   \left(1+\frac{R_{36}}{\ell^{36}}\right)^{-2}=1-p,\notag \\
 \Pt(72,b_0)&\ge
   \left(U^{72}+R_{72}\right)^{-2}=1-q.
 \label{cert:eq:base-purities}
\end{align}
For $\Ps(36,b_0)$, we bound the numerator from below by
$Z_{72}(b_0)\ge\lambda_0^{72}$, bound $Z_{36}(b_0)\le|\lambda_0|^{36}+R_{36}$,
and use $|\lambda_0|>\ell$. For $\Pt(72,b_0)$, we bound the numerator from
below by $Z_{72}(2b_0)>1$ from Eq.~\eqref{cert:eq:trial-consequence} and
bound the denominator by $Z_{72}(b_0)=\lambda_0^{72}+\sum_{i\ne0}\lambda_i^{72}
\le U^{72}+R_{72}$.

Because $k/12\in\{3,6\}$ is an integer, Eqs.~\eqref{cert:eq:dominance-check}--\eqref{cert:eq:base-purities}
involve only rational arithmetic. Exact evaluation verifies
\begin{equation}
 p\le\frac{47916}{10^6},\qquad
 q\le\frac{181521}{10^6},
 \label{cert:eq:scalar-checks}
\end{equation}
which completes the proof of Eq.~\eqref{cert:eq:base}. Numerically,
$p\simeq0.0479157710062$ and $q\simeq0.181520168033$.

\subsection{Temperature lifting and the improved uniform bound}
\label{cert:sec:lifting}

Because the gap bound from the doubling argument scales inversely with the
seed inverse temperature, certifying smaller defects at a moderate inverse
temperature and stopping after a finite prefix of doubling steps gives a
much larger gap constant. We now explain the temperature-lifting method,
which uses the same trace table at $b_0=49/4$ and the $120$-site trial
bound $Z_{120}(\beta)>1$ to certify a seed at $(n_0,b)=(60,rb_0)$ with
$r\ge1$. Combining that lifted seed with a tighter unlifted bound at
$(36,b_0)$ covers all even lengths $L\ge36$.

\subsubsection{Lifted moments and a dominant spatial weight}
\label{cert:sec:lifted-moments}

We keep $m_J,m_N$ from Eq.~\eqref{cert:eq:masses} and sharpen the filter
caps~\eqref{cert:eq:caps} to the exact rationals
\begin{equation}
 U=1.001350644012,\qquad V=0.871255704653,
 \label{cert:eq:tight-caps}
\end{equation}
verified by the same exterior-ray expansions of $F_J$ and $F_N$.
Since $U^{12}<m_J<2U^{12}$ and $V^{12}<m_N<2V^{12}$, Eq.~\eqref{cert:eq:capped}
bounds the full sector moments at any even length $k\ge12$ by
\begin{align}
 \mathcal J_k&=U^k+(m_J-U^{12})^{k/12},\notag \\
 \mathcal N_k&=V^k+(m_N-V^{12})^{k/12},\notag \\
 M_k&=\mathcal J_k+3\mathcal N_k.
 \label{cert:eq:moment-upper}
\end{align}
Because $N_k\ge0$ at even $k$, we have $Z_k(b_0)=J_k+3N_k\le M_k$ and
$Z_k(b_0,P)=J_k-N_k\le\mathcal J_k$.

Now fix a rational factor $r\ge1$, an even length $12\le k\le60$, and the
lifted inverse temperature $b=rb_0$. Both $H_k$ and $H_k(P)$ are
self-adjoint, so their physical Boltzmann weights $z_i>0$ at $b_0$ satisfy
$\sum_i z_i^r\le(\sum_i z_i)^r$. This transfers the upper bounds at $b_0$
to $b$:
\begin{equation}
 Z_k(b)\le D:=M_k^r,\qquad
 Z_k(b,P)\le E:=\mathcal J_k^r.
 \label{cert:eq:temperature-lift}
\end{equation}
In rational arithmetic we replace $D$ and $E$ by outward rational upper
bounds satisfying $E\le D$.

At the lifted inverse temperature $b$, let $K_b$ be the one-site spatial
transfer operator and define the $k$th-power weights
\begin{equation}
 w_i=|\lambda_i(b)|^k,\qquad h=\frac{60}{k},\qquad
 \tau=\frac{D+3E}{4}.
 \label{cert:eq:spatial-weights}
\end{equation}
Since $k$ is even, $w_i=\lambda_i(b)^k\ge0$. Thus $\sum_i w_i=Z_k(b)\le D$,
the $J$-sector sum satisfies $\sum_{\lambda_i\in J}w_i=[Z_k(b)+3Z_k(b,P)]/4\le\tau$,
and the $120$-site trial bound~\eqref{cert:eq:trial-consequence} gives
\begin{equation}
 \sum_iw_i^{2h}=Z_{120}(b)>1.
 \label{cert:eq:lifted-trial}
\end{equation}
If a rational threshold $t>0$ satisfies
\begin{equation}
 Dt^{2h-1}<1,\qquad 3t>D,
 \label{cert:eq:weight-forcing}
\end{equation}
then at least one weight must satisfy $w_0>t$; otherwise
$\sum_i w_i^{2h}\le t^{2h-1}\sum_i w_i\le Dt^{2h-1}<1$, contradicting
Eq.~\eqref{cert:eq:lifted-trial}. Moreover, $3t>D$ forces $w_0$ to lie in
$J$, because any weight in $N$ is triply degenerate and would contribute
more than $3t>D$ to $\sum_i w_i$.

Removing $w_0$ leaves residual $J$ mass $x\in[0,\tau-t]$, while each of
the three identical $N$ sectors has mass at most $(D-t-x)/3$. For any
$s\ge1$, the residual moment $\sum_{i\ne0}w_i^s$ is at most
$x^s+3[(D-t-x)/3]^s$. By convexity on $x\in[0,\tau-t]$, this sum attains
its maximum at one of the two endpoints:
\begin{equation}
 \mathcal R_s(t)=
 \max\left\{3\left(\frac{D-t}{3}\right)^s,
 (\tau-t)^s+3\left(\frac{D-\tau}{3}\right)^s\right\}.
 \label{cert:eq:lifted-residual}
\end{equation}
All choices below satisfy $0<t\le\tau\le D$.

Once $w_0>t$ is established, any larger rational $t'$ with
\begin{equation}
 (t')^{2h}+\mathcal R_{2h}(t)<1
 \label{cert:eq:refinement}
\end{equation}
implies $w_0>t'$, since $w_0\le t'$ would give
$Z_{120}(b)=w_0^{2h}+\sum_{i\ne0}w_i^{2h}\le(t')^{2h}+\mathcal R_{2h}(t)<1$.
Iterating Eq.~\eqref{cert:eq:refinement} tightens the lower bound $t$ on
$w_0$. Because $w_0^h=|\lambda_0(b)|^{60}$ and
$\sum_{i\ne0}|\lambda_i(b)|^{60}\le\mathcal R_h(t)$, we obtain the lifted
spatial purity bound at $n=60$:
\begin{equation}
 \Ps(60,b)\ge
 \left(1+\frac{\mathcal R_h(t)}{t^h}\right)^{-2}.
 \label{cert:eq:lifted-spatial}
\end{equation}

\subsubsection{The physical thermal residual}
\label{cert:sec:thermal-residual}

To bound $\Pt(120,b)$ at $b=rb_0$, we work directly with the physical
Boltzmann weights $z_i=e^{-b_0\widehat E_i(120)}$ of $\widehat H_{120}$
rather than the spatial eigenvalues of $K_b$. Ordering $z_0\ge z_1\ge\cdots$,
we have $z_0>1$ from the $120$-site trial state and
$\sum_i z_i=Z_{120}(b_0)\le M_{120}$ from Eq.~\eqref{cert:eq:moment-upper}.
Hence
\begin{equation}
 \sum_{i>0}\frac{z_i}{z_0}\le M_{120}-1,\qquad
 \sum_{i>0}\left(\frac{z_i}{z_0}\right)^r
 \le(M_{120}-1)^r.
 \label{cert:eq:boltzmann-residual}
\end{equation}
At $b=rb_0$, the physical partition functions are $Z_{120}(b)=\sum_i z_i^r$
and $Z_{120}(2b)=\sum_i z_i^{2r}\ge z_0^{2r}$. Dividing numerator and
denominator by $z_0^{2r}$ gives the lifted thermal purity bound:
\begin{equation}
 \Pt(120,rb_0)\ge
 \left[1+(M_{120}-1)^r\right]^{-2}.
 \label{cert:eq:lifted-thermal}
\end{equation}
The transfer identity $\Tr K_\beta^L=Z_L(\beta)$ equates the full traces
at each $\beta$; it does not identify individual spatial eigenvalues with
physical Boltzmann weights.

\subsubsection{Finite-prefix coverage and the doubling tail}
\label{cert:sec:coverage}

We now convert the seed bounds at $(n_0,b_0)$ into a uniform gap bound by
combining a finite prefix of rational doubling steps with an analytic tail
estimate.

\begin{proposition}[Finite-prefix coverage]
\label{cert:prop:prefix}
Let $n_0$ be even and $b_0>0$. Suppose $0\le p_j,q_j<1$ bound respectively
$1-\Ps(n_j,b_j)$ and $1-\Pt(2n_j,b_j)$ at
$(n_j,b_j)=(2^jn_0,2^jb_0)$ for $0\le j\le K$.
Put $A_j=(1-p_j)^2(1-q_j)$ and assume $A_j>1/2$ for $j<K$.
Let $u=\max(p_K,q_K)$ and choose $C$ such that
\begin{align}
 &0<u<1/6,\qquad Cu<1,\qquad C(1-3u)>3,\notag \\
 G(u):&=\frac12\left((1-u)^{-1}+(1-u)^{-2}
                         +(1-u)^{-3}\right)^2\le C.
 \label{cert:eq:tail-checks}
\end{align}
If $g$ is strictly smaller than every member of
\begin{equation}
 \left\{\frac{\log(A_j/(1-A_j))}{b_j}:0\le j<K\right\}
 \cup\left\{\frac{\log(1/(Cu))}{b_K}\right\},
 \label{cert:eq:prefix-minimum}
\end{equation}
then every even $L\ge n_0$ has a unique ground state and
$\Delta_L>g$. An entry with $A_j=1$ is interpreted as $+\infty$.
\end{proposition}

\begin{proof}
For even $n_j\le L\le2n_j$ with $0\le j<K$, spatial power-sum
interpolation~\eqref{sm:eq:interpolation} gives
$\Pt(L,b_j)\ge A_j^{L/(2n_j)}\ge A_j>1/2$.
The largest normalized Boltzmann weight $a_0$ satisfies $a_0\ge\Pt(L,b_j)>1/2$,
so the ground state is unique, and the next weight $a_1\le1-a_0$ obeys
\begin{equation}
 e^{-b_j\Delta_L}=\frac{a_1}{a_0}
 \le\frac{1-a_0}{a_0}\le\frac{1-A_j}{A_j},
 \label{cert:eq:interval-gap}
\end{equation}
which yields $\Delta_L\ge\log(A_j/(1-A_j))/b_j>g$.
For $j\ge K$, since $f(1-(1-x)^3)=x^2G(x)$ and $G$ is increasing on
$[0,u]$, Eq.~\eqref{cert:eq:tail-checks} implies that $u_m=C^{-1}(Cu)^{2^{m-K}}$
bounds both defects at $(n_m,b_m)$ by induction. Interpolating on
$n_m\le L\le2n_m$ gives $\Pt(L,b_m)\ge1-3u_m>1/2$ and
\[
 e^{-b_m\Delta_L}\le\frac{3u_m}{1-3u_m}
 \le\frac{3}{C(1-3u)}(Cu)^{2^{m-K}}<(Cu)^{2^{m-K}},
\]
so $\Delta_L>\log(1/(Cu))/b_K>g$. Together, the intervals $[n_j,2n_j]$
for $0\le j<K$ and the tail $L\ge n_K$ cover every even $L\ge n_0$.
\end{proof}

We generate the prefix $(p_j,q_j)$ from the asymmetric recurrence with
$f(x)=x^2/[2(1-x)^2]$:
\begin{align}
 p_{j+1}&=\left\lceil
 f(1-(1-p_j)^2(1-q_j))\right\rceil_{\delta},\notag \\
 q_{j+1}&=\left\lceil
 f(1-(1-q_j)^2(1-p_{j+1}))\right\rceil_{\delta},
 \label{cert:eq:recurrence}
\end{align}
where $\lceil x\rceil_\delta=\delta\lceil x/\delta\rceil$.
For the lifted seeds at $(60,rb_0)$, we take $K=10$ steps with
$\delta=10^{-2048}$ and apply Proposition~\ref{cert:prop:prefix}.

\begin{table}[t]
 \caption{Temperature-lifting certificates at $(60,rb_0)$ with $b_0=49/4$.
 Defect bounds at $j=0$ are rounded upward; gap bounds hold for every even
 $L\ge60$, and for every even $L\ge36$ once combined with the unlifted
 bridge on $36\le L\le72$.}
 \label{cert:tab:lifted}
 \begin{ruledtabular}
 \begin{tabular}{ccccc}
 $k$ & $r$ & Spatial defect & Thermal defect & $\Delta_L>$\\
 \hline
 30 & $2$    & 0.0307608679 & 0.0590740709 & 0.070450002\\
 28 & $17/8$ & 0.0305302052 & 0.0479556776 & 0.072901771
 \end{tabular}
 \end{ruledtabular}
\end{table}

Table~\ref{cert:tab:lifted} lists two parameter choices. The integer choice
$(k,r)=(30,2)$ has $h=2$, so its spatial residual~\eqref{cert:eq:lifted-residual}
uses only second and fourth powers and already gives $\Delta_L>0.07045$
for even $L\ge60$. The fractional choice $(k,r)=(28,17/8)$ gives
$\Delta_L>0.072901771$ for even $L\ge60$.

To cover the remaining even lengths $36\le L<60$, we evaluate the unlifted
base-seed formulas~\eqref{cert:eq:base-residual}--\eqref{cert:eq:base-purities}
at $(36,b_0)$ using the sharper caps~\eqref{cert:eq:tight-caps} and
$\ell=0.999997937629$, validated by the capped-mass inequality
$\mathcal C(m_J,\ell^{12},6)+3\mathcal C(m_N,V^{12},6)<1$.
At $j=0$, the resulting factor $A_0=(1-p_0)^2(1-q_0)$ satisfies
$\log(A_0/(1-A_0))/b_0>0.09$, which by Eq.~\eqref{cert:eq:interval-gap}
bounds $\Delta_L>0.09$ on every even $36\le L\le72$. Combining this bridge
on $[36,72]$ with the lifted bound on $[60,\infty)$ proves
Eq.~\eqref{cert:eq:improved} for every even $L\ge36$.

\subsection{Reproducibility and arithmetic assumptions}
\label{cert:sec:reproduce}

The \href{https://github.com/LionSR/CertifiedGaps}{repository} provides
the stored certificates and standalone verifiers for the finite inputs and
scalar inequalities above.

To keep every scalar check in exact rational arithmetic, we enclose
rational powers and logarithms by rational endpoints. For $x\ge0$ and a
rational exponent $a/b\ge0$ with $b>0$, nonnegative rational bounds $l,u$
satisfy $l\le x^{a/b}\le u$ if and only if the integer powers satisfy
$l^b\le x^a\le u^b$. To enclose $\log x$ for $x>0$, we write $x=2^ky$ with
$k\in\mathbb Z$ and $1\le y<2$, and set $z=(y-1)/(y+1)\in[0,1/3)$. The
series
\begin{align}
 \log y&=2\sum_{j=0}^{M-1}\frac{z^{2j+1}}{2j+1}+E_M,\notag \\
 0\le E_M&\le
 \frac{2z^{2M+1}}{(2M+1)(1-z^2)}
 \label{cert:eq:log-enclosure}
\end{align}
at $M=40$ encloses $\log y$ and $\log2$ in rational intervals, which we
round outward. Floating-point calculations are used only to select
parameters, never in the certificate checks.

Our trusted base consists of exact integer and rational arithmetic, the
explicit Poisson and Horner error bounds for the short-ring traces, and
the analytical transfer, interpolation, and doubling inequalities proved
in Sections~\ref{sm:sec:transfer}--\ref{sm:sec:gap} and above. Together
with the finite trace and trial inputs of Ref.~\cite{OpenAI2026Periodic},
these checks establish Eq.~\eqref{cert:eq:improved} for the periodic
spin-1 chain~\eqref{cert:eq:hamiltonian}.
